\section{Introduction}\label{sec:intro}

Viability theory asks from which states of a constrained control system at least one control can keep every future state inside a constraint set \citep{aubin1991}. For sustainability problems the constraints are floors---stock levels, service thresholds, safety margins---and viability is the formal counterpart of maintaining a development path rather than optimising it \citep{bene2001,delara2008,doyen2012,doyen2020}. Under \emph{perfect measurement} the theory is complete in the relevant sense: the viability kernel exists as the largest viable subset of the constraint set, Nagumo-type tangency conditions characterise it, and approximation schemes compute it \citep{saintpierre1994,aubin2011}.

Sustainability governance, however, operates under \emph{incomplete observation}. Stocks are assessed at discrete review intervals; indicators are coarser than the state; some stocks are not observed at all; and the information that does arrive may arrive late. The signal that would reveal an erosion of the productive base is often precisely the one that is coarse, postponed, or absent, because a base can be drawn down while the measured output is held up. Incomplete observation is therefore structural to the certification problem, not a practical nuisance.

The viability question under incomplete observation splits into two directions. The sufficiency direction has a canonical literature. \citet{veliov1993} gives conditions for the existence of an output-feedback regulation map under incomplete and inexact measurement, reducing under perfect measurement to the classical viability condition of \citet{haddad1981}; the estimation-tube programme absorbs imperfect measurement into set-valued dynamics on a space of sets of compatible states \citep{kurzhanski1997,cardaliaguet2007}; barrier certificates certify safety from a scalar function \citep{prajna2004,prajna2007}, with converse and hybrid characterisations \citep{prajna2005,maghenem2019}; and Hamilton--Jacobi reachability computes backward reachable tubes as sublevel sets of viscosity solutions \citep{mitchell2005,margellos2011,bansal2017}. All of these are, in direction, sufficiency machines: they compute or certify the states from which some control preserves safety.

What that literature does not supply is the converse instrument: a systematic way of \emph{certifying that no observation-based policy is viable}, and of saying which feature of the observation design---its timing, coarseness, aggregation or bias---is responsible. This paper develops that instrument. An \emph{obstruction certificate} is a checkable witness that a prescribed class of observation-based policies fails, not because a particular policy is bad, but because the information structure leaves no room for any policy. Its negation is a necessary condition for observation-based viability. In the polyhedral and finite cases the certificate is a finite object---a Farkas pair, a short list of incompatible states, an obstruction tree---and in the analytic cases it is an explicit adverse trajectory together with a quantitative bound.

\paragraph{The organising idea.} An observation-based regulator cannot tailor its action to the true state; it must select a single response---a held command, a blind-window control, or a safe/unsafe verdict---for an entire information state. Observation-based viability is therefore the existence of one response that is simultaneously admissible, safe until the next observation, and recursively viable, for every state compatible with the record (the \emph{selector principle}, Proposition~\ref{prop:selector}). Each certificate of the paper proves that one of these response correspondences is empty, and the certificates are nested in a ladder. This single viewpoint turns a list of nonviability conditions into a calculus: the mechanism that fires identifies the correspondence that failed, and the correspondence that failed identifies the design change that could restore it.

\paragraph{Contributions.} The paper establishes the following, each stated with the policy class and the hypotheses under which it holds.
\begin{enumerate}[leftmargin=2em,itemsep=2pt]
\item \emph{Exactness in finite systems} (Section~\ref{sec:finite}). Backward recursion over set-valued beliefs is sound and complete on every finite horizon (Theorem~\ref{thm:finite-horizon}), and one step ahead the common-action test is necessary and sufficient (Theorem~\ref{thm:onestep}); nonviability is witnessed by a finite obstruction tree. This fragment is machine-checked in Lean~4.
\item \emph{The common-action obstruction} (Theorem~\ref{thm:common-action}). At a post-observation information state of a held-action review class, emptiness of the common safe-action set---under a compatible adverse-path hypothesis---excludes every policy, even when each compatible state is viable under full information. For polyhedral data the witness is a Farkas pair; for control-affine dynamics with convex control sets a Helly-type argument bounds the witness by $m+1$ compatible states (Proposition~\ref{prop:helly}). The rungs of the obstruction ladder (Proposition~\ref{prop:ladder}) are the admissibility, safe-action and review-tube correspondences.
\item \emph{The delayed-information obstruction} (Theorem~\ref{thm:delayed}). When the first branch-discriminating observation arrives at time $\Tobs$ and an adverse realisation drives a constraint margin down at rate $\varepsilon$, the explicit bound $\inf q/\varepsilon<\Tobs$ excludes every policy of the blind-window class; the bound has a sharp threshold form (Remark~\ref{rem:sigma}), and over polytope-declared blind-window classes the threshold is exactly computable as a linear programme whose Farkas dual is the certificate (Theorem~\ref{thm:lp}). The class declaration is load-bearing: a whole-window hold class can have an unbounded survival gap relative to its per-step polytope relaxation (Example~\ref{ex:relax-gap}).
\item \emph{Dynamic and recourse obstructions} (Theorem~\ref{thm:exit}, Proposition~\ref{prop:recourse}). A finite-time exit certificate of comparison type closes the viability question before any observation argument, and an oracle-recourse certificate for affine systems detects the complementary mode in which a blind window is survivable but no post-observation authority can recover every branch; a planar braking instance has an exact delay threshold $\tau^*=7/50$ and is genuinely three-way.
\item \emph{Certification limits} (Section~\ref{sec:certification}). An exact observation-only safety certifier exists if and only if safe-set membership is constant on the observation fibres (Proposition~\ref{prop:fibre}); the certainly-safe set is the sound relaxation (Corollary~\ref{cor:certainly-safe}); and a certainty-equivalence controller applied to a biased but injective observation can have an empty kernel while a corrected controller is viable (Proposition~\ref{prop:ce-trap}).
\item \emph{The limits of completeness} (Section~\ref{sec:completeness}). Static observation reduces exactly to open-loop robust viability (Theorem~\ref{thm:static}); the two-phase (blind window, then reveal) problem decomposes exactly (Proposition~\ref{prop:decomposition}); and no certificate pair that reads only the window sub-model can be complete (Proposition~\ref{prop:window-nogo}). The residual dynamic gap is stated as an open problem with a definite target.
\end{enumerate}
The certificates are instantiated on a symbolic resource model and audited, cell by cell, against the exact recursion (Section~\ref{sec:casestudy}); their monitoring and institutional consequences---timing, coarseness, aggregation and bias---are drawn in Section~\ref{sec:discussion}.

\paragraph{Related work.} The nearest neighbour of the question posed here is \citet{doyen2000}, who characterises \emph{guaranteed viability domains} under output feedback: a closed set $K$ is such a domain when there is a Lipschitz selection $u(\cdot)$ of the control constraint map with $K$ invariant under $\dot x\in\{f(x,u(h(x,w)),w):w\in W(x)\}$, and the characterisation is an equivalence whose constructive content is a feedback synthesis by Steiner selection. The present work differs on three axes visible in that statement: the policy class (memoryless Lipschitz functions of the current output there; measurable nonanticipative functions of the whole record here), the property (exact invariance of a closed set there; viability of information states here), and the direction---the obstruction side, which \citet{doyen2000} records as undeveloped (``no general viability result is available in this differential game context with imperfect and/or partial information''). On the failure side under full information, \citet{aubin2001} characterises the complement of the kernel together with capture basins, and \citet{aubin2002} supply the bilateral-fixed-point calculus; the certificates below are one-sided, observation-constrained analogues of that programme. Robust viability is \citet{aubinfrankowska1990} and \citet{frankowska1989}; viability with observable parameters is \citet{quincampoix1994}. The set-membership information state used throughout is that of \citet{kurzhanski1997}; the finite belief recursion is the set-valued analogue of the belief-state recursion of partially observable Markov processes \citep{astrom1965,smallwood1973}, whose qualitative safety variant is \citet{chatterjee2009}, and whose complexity is \citet{papadimitriou1987}. Hamilton--Jacobi reachability \citep{mitchell2005,margellos2011,bansal2017}, resting on viscosity-solution theory \citep{crandall1992}, indexes its value by a physical state, whereas this paper tracks sets of compatible states; its learning-based extensions \citep{bansal2021} purchase scalability with approximation. Shield-based runtime enforcement for finite stochastic models \citep{alshiekh2018} restricts an arbitrary policy to precomputed safe actions and is a sufficiency device in the present sense. A companion robust-viability application to a fishery limit reference point is \citet{abaee2026fish}, and the minimax-dual development of the common-action witnesses is a separate companion manuscript \citep{abaee2026minimax}. To the author's knowledge the common-action and delayed-information obstructions and the fibre criterion have not been stated in this form; the elementary facts on which they rest---quantifier commutation, Dini comparison, Farkas' lemma and Helly's theorem---are classical.

\paragraph{Organisation.} Section~\ref{sec:setting} fixes the setting, consolidates the standing conventions, and states the selector principle. Section~\ref{sec:finite} treats finite systems, where the calculus is exact. Section~\ref{sec:calculus} develops the continuous-time certificates. Section~\ref{sec:certification} treats certification limits, Section~\ref{sec:completeness} the limits of completeness, Section~\ref{sec:casestudy} the worked model and coverage audit, and Section~\ref{sec:discussion} the position of the calculus and its design consequences. Short proofs are given in place; complete proofs of the longer statements, the worked certificates, the probabilistic counterpart, the expanded review of the sufficiency literature, and the verification record are in the Supplementary Material (cited as S1--S9). Table~S9.1 maps every result of this edition to its predecessor.
